% MG06. In exc:paley-codigo, after defining C and chi, before the
% quadratic-sum sentence and reduction modulo three.
The quadratic identity in this chart follows from a finite
correlation. The order \((\infty,0,1,2,3,4)\) is retained,
together with the declared identification of the six positions
and the character \(\chi\).

\begin{lemma}[Correlation of the quadratic character on five positions]
\label{mg06:correlacion}
For \(a,b\in\mathbb F_5\) with \(a\ne b\),
\begin{equation}
 \sum_{t\in\mathbb F_5}\chi(t-a)\chi(t-b)=-1.
 \label{mg06:identidad-correlacion}
\end{equation}
\end{lemma}
\begin{proof}
The substitution \(u=(t-a)/(b-a)\) is bijective because \(b-a\)
is invertible in \(\mathbb F_5\). Furthermore,
\[
 (t-a)(t-b)=(b-a)^2u(u-1),\qquad
 \chi((b-a)^2)=1.
\]
Multiplicativity of \(\chi\), including its extension by zero,
reduces the sum to \(\sum_u\chi(u(u-1))\). The five evaluations are
\begin{equation}
\begin{array}{c|rrrrr}
 u&0&1&2&3&4\\\hline
 u(u-1)\pmod5&0&0&2&1&2\\
 \chi(u(u-1))&0&0&-1&1&-1
\end{array}
\label{mg06:cinco-evaluaciones}
\end{equation}
Their sum is \(-1\), proving
\eqref{mg06:identidad-correlacion}.
\end{proof}

\Needspace{8\baselineskip} % REV03: keep the destination with the proposition heading.
\begin{proposition}[Orthogonality of the six-position chart]
\label{mg06:conferencia}
The matrix \(C\) satisfies
\begin{equation}
 C=C^{\mathsf T},\qquad CC^{\mathsf T}=5I_6,
 \qquad C^2=5I_6.
 \label{mg06:identidad-conferencia}
\end{equation}
\end{proposition}
\begin{proof}
Since \(\chi(-1)=\chi(4)=1\), one has
\(\chi(b-a)=\chi(a-b)\); entries involving \(\infty\) are
also symmetric. Each row contains one zero entry and five entries
of absolute value one, so its squared norm is five. For two distinct
finite rows, the inner product receives \(+1\) from the
\(\infty\) coordinate and \(-1\) from the five finite coordinates
by \eqref{mg06:identidad-correlacion}; their sum is zero.
The inner product of the infinite row with a finite row is
\(\sum_{t\in\mathbb F_5}\chi(t)=0+1-1-1+1=0\).
This evaluates every entry of \(CC^{\mathsf T}\).
Symmetry converts this identity into \(C^2=5I_6\).
\end{proof}

The argument retains the chart order to be reduced modulo three.
The quadratic identity and the subsequent graph code therefore use
the same coordinates, with their incidences and supports.
